% Part: first-order-logic
% Chapter: sequent-calculus
% Section: derivations

\documentclass[../../../include/open-logic-section]{subfiles}

\begin{document}

\iftag{FOL}
      {\olfileid{fol}{seq}{der}}
      {\olfileid{pl}{seq}{der}}

\olsection{\usetoken{P}{derivation}}

\begin{explain}
موږ وويل چې لومړنے سېکوېنټ څنګه وي او د استنتاج قاعدې مو ورکړې۔
په سېکوېنټ حساب کښې !!^{derivation}s له دغو څخه په استقرايي توګه
جوړېږي: هر !!{derivation} يا په يوازې ځان يو لومړنے سېکوېنټ وي،
او يا له يوه يا دوو !!{derivation}s او ورپسې يوه استنتاج څخه جوړ وي۔
\end{explain}

\begin{defn}[$\Log{LK}$ !!{derivation}]
د يوه سېکوېنټ~$S$ \emph{$\Log{LK}$-!!{derivation}} د سېکوېنټونو
داسې متناهي ونه ده چې لاندې شرطونه پوره کوي:
\begin{enumerate}
\item د ونې تر ټولو بره سېکوېنټونه لومړني سېکوېنټونه دي۔
\item د ونې تر ټولو لاندې سېکوېنټ~$S$ دے۔
\item له~$S$ پرته د ونې هر سېکوېنټ د استنتاج د يوې قاعدې د سمې
  کارونې مقدمه ده، او د هغې کارونې پايله په ونه کښې د هماغه
  سېکوېنټ مخامخ لاندې ولاړه وي۔
\end{enumerate}
بيا وايو چې $S$ د !!{derivation} \emph{وروستنے سېکوېنټ} دے او
$S$ \emph{په $\Log{LK}$ کښې !!{derivable}} دے (يا
$\Log{LK}$-!!{derivable} دے)۔
\end{defn}

\begin{ex}
هر لومړنے سېکوېنټ، مثلاً $!C \Sequent !C$، !!a{derivation} دے۔
له دې څخه د يوې بلې قاعدې، مثلاً د $\LeftR{\Weakening}$ قاعدې، په
کارولو نوے !!{derivation} ترلاسه کولے شو:
\begin{prooftree}
\Axiom$ \Gamma \fCenter \Delta $
\RightLabel{\LeftR{\Weakening}}
\UnaryInf$ !A, \Gamma \fCenter \Delta$
\end{prooftree}
خو قاعده عمومي مراد لري: په قاعده کښې $!A$ د هرې !!{sentence}،
مثلاً د~$!D$، په ځاے راوستلے شو۔ که مقدمه زموږ له لومړني سېکوېنټ
$!C \Sequent !C$ سره سمون خوري، نو $\Gamma$ او $\Delta$ دواړه
يوازې~$!C$ دي، او پايله به بيا $!D, !C \Sequent !C$ وي۔ نو لاندې
!!a{derivation} دے:
\begin{prooftree}
\Axiom$ !C \fCenter !C $
\RightLabel{\LeftR{\Weakening}}
\UnaryInf$ !D, !C \fCenter !C$
\end{prooftree}
اوس يوه بله قاعده، مثلاً $\LeftR{\Exchange}$، کارولے شو چې په کيڼ
لوري کښې د دوو !!{sentence}s د ځاے بدلولو اجازه راکوي۔ نو لاندې هم
سم !!{derivation} دے:
\begin{prooftree}
\Axiom$ !C \fCenter !C $
\RightLabel{\LeftR{\Weakening}}
\UnaryInf$ !D, !C \fCenter !C$
\RightLabel{\LeftR{\Exchange}}
\UnaryInf$ !C, !D \fCenter !C$
\end{prooftree}
د دې قاعدې په دې کارونه کښې، چې داسې ورکړل شوې وه:
\begin{prooftree}
\Axiom$ \Gamma, !A, !B, \Pi \fCenter \Delta $
\RightLabel{\LeftR{\Exchange}}
\UnaryInf$ \Gamma, !B, !A, \Pi \fCenter \Delta,$
\end{prooftree}
$\Gamma$ او $\Pi$ دواړه تش وو، $\Delta$ هم $!C$ دے، او د $!A$ او
$!B$ ونډې په همدې ترتيب $!D$ او~$!C$ لوبوي۔ نژدې په همدې طريقه
وينو چې
\begin{prooftree}
\Axiom$ !D \fCenter !D $
\RightLabel{\LeftR{\Weakening}}
\UnaryInf$ !C, !D \fCenter !D$
\end{prooftree}
هم !!a{derivation} دے۔ اوس دا دواړه !!{derivation}s اخلو او د
$\RightR{\land}$ په کارولو ئې يو ځاے کوو۔ هغه قاعده داسې وه:
\begin{prooftree}
\Axiom$\Gamma \fCenter \Delta, !A$
\Axiom$ \Gamma \fCenter \Delta, !B$
\RightLabel{\RightR{\land}}
\BinaryInf$ \Gamma \fCenter \Delta, !A \land !B $
\end{prooftree}
زموږ په حالت کښې مقدمې بايد د هغو !!{derivation}s له وروستيو
سېکوېنټونو سره سمون وخوري چې په مقدمو پاے ته رسېږي۔ يعنې $\Gamma$
د $!C, !D$، $\Delta$ تشه، $!A$ د $!C$ او $!B$ د $!D$ سره برابره
ده۔ نو که استنتاج سم وي، پايله $!C, !D \Sequent !C
\land !D$ ده۔
\begin{prooftree}
\Axiom$ !C \fCenter !C $
\RightLabel{\LeftR{\Weakening}}
\UnaryInf$ !D, !C \fCenter !C$
\RightLabel{\LeftR{\Exchange}}
\UnaryInf$ !C, !D \fCenter !C$
\Axiom$ !D \fCenter !D $
\RightLabel{\LeftR{\Weakening}}
\UnaryInf$ !C, !D \fCenter !D$
\RightLabel{\RightR{\land}}
\BinaryInf$ !C, !D \fCenter !C \land !D $
\end{prooftree}
البته مقدمې سرچپه کولے هم شو؛ بيا $!A$ به $!D$ او $!B$ به~$!C$ وي۔
\begin{prooftree}
\Axiom$ !D \fCenter !D $
\RightLabel{\LeftR{\Weakening}}
\UnaryInf$ !C, !D \fCenter !D$
\Axiom$ !C \fCenter !C $
\RightLabel{\LeftR{\Weakening}}
\UnaryInf$ !D, !C \fCenter !C$
\RightLabel{\LeftR{\Exchange}}
\UnaryInf$ !C, !D \fCenter !C$
\RightLabel{\RightR{\land}}
\BinaryInf$ !C, !D \fCenter !D \land !C $
\end{prooftree}
\end{ex}

\end{document}
